\textbf{COLUNI 2025 -- Questão 20}\par\medskip “Em março de 2024, o custo da cesta básica da cidade de São Paulo foi o maior entre as 17 capitais onde o Departamento Intersindical de Estatística e Estudos Socioeconômicos (DIEESE) realiza mensalmente a Pesquisa Nacional da Cesta Básica de Alimentos, chegando a R\$ 813,26, alta de 0,60\% em relação a fevereiro de 2024. Em comparação com março de 2023, a cesta subiu 3,97\% [...]. Considerando-se a determinação constitucional que estabelece que o salário-mínimo deve ser suficiente para suprir as despesas de um trabalhador e da família dele com alimentação, moradia, saúde, educação, vestuário, higiene, transporte, lazer e previdência, [...] o DIEESE estima que, nesse caso, [...] o salário-mínimo necessário para a manutenção de uma família de quatro pessoas deveria ter sido 4,84 vezes o mínimo reajustado em R\$ 1.412,00”.

\noindent Disponível em: https://www.dieese.org.br/analisecestabasica/2024/202403cestabasica.pdf. Acesso em: 16 jul. 2024.

\par

Com base nos dados acima, é CORRETO afirmar que:
\begin{enumerate}
\item O preço da cesta básica, em fevereiro de 2024, foi de, aproximadamente, R\$ 798,32.
\item O salário-mínimo necessário, em março de 2024, deveria ser R\$ 5.834,08.
\item O salário-mínimo necessário, em março de 2024, deveria ser R\$ 6.804,08.
\item O preço da cesta básica, em março de 2023, foi de, aproximadamente, R\$ 782, 21.
\end{enumerate}
